%% kapa: exam=1 round=25 year=2014 period=1 form=A subject=economics subject_ko=경제원론 seq=13 number=53
%% source: Q-Net 감정평가사 기출문제 (cst00309) artlSeq=5101207 file=extracted/1008/5101207/2014년 제25회 감정평가사 1차시험 기출문제/1교시(A형).hwp
%% conversion: kordoc 4.24.51 -> tools/gen1.py (auto); review=proofread-sonnet; flags=none
\begin{question}{53}
과수원주인인 甲과 양봉업자인 乙이 인근 지역에서 경제활동을 하고 있는데, 甲이 과실나무를 더 많이 심자 乙의 꿀 생산이 증가하고, 乙이 꿀벌의 수를 증가시키자 과수원 수확이 늘어나는 것을 확인할 수 있었다. 甲과 乙에게 발생하는 외부성에 관한 설명으로 옳은 것을 모두 고른 것은?
\begin{framedtext}
ㄱ. 甲과 乙은 각각 서로에게 양의 외부성을 주게 된다. \par
ㄴ. 거래비용이 존재하지 않을 때, 甲과 乙 간의 거래에 의해 사회적 최적 생산량을 합의해낼 수 있다. \par
ㄷ. 甲과 乙 사이에 서로 양의 외부성을 주고받는 경우이므로, 시장실패에 대한 교정은 불필요하다. \par
ㄹ. 甲이 양봉장을 인수함으로써 사회적 최적생산량을 달성할 수 있다. \par
\end{framedtext}
\begin{choices}
  \item ㄱ, ㄴ
  \item ㄴ, ㄷ
  \item ㄷ, ㄹ
  \item ㄱ, ㄴ, ㄹ
  \item ㄱ, ㄷ, ㄹ
\end{choices}
\end{question}
\answer{4}
